\begin{helper}
Fix \(\eta>0\), let \(\beta=1/2\), and let
\[
m(n)=M(n)=\noderef{TwoBiteNaturalReducedVertexCount}(n).
\]
Put \(p(n)=\noderef{TwoBiteEdgeProbability}(\beta,n)\).  Under
\(\noderef{TwoBiteEventProbability}(n,m(n),p(n))\), the simultaneous event
\[
\forall x\in \noderef{TwoBiteBaseVertex}(m(n)),\qquad
\noderef{WithinRelativeError}\!\left(
  \eta,\ p(n)n,\ |\noderef{TwoBiteLiftedNeighborhood}(\omega,x)|
\right)
\]
holds with high probability in the sense of \(\noderef{WithHighProbability}\).
\end{helper}
\begin{proof}
Write \(L=\log n\), \(M=M(n)\), and
\[
p=p(n)=\frac12\sqrt{\frac{L}{n}}.
\]
By \(\noderef{TwoBiteNaturalReducedVertexCount}\), together with
\(\noderef{TwoBiteReducedVertexCount}\) and \(\noderef{TwoBiteStretch}\), for
all sufficiently large \(n\),
\[
M=\left\lfloor\frac{n}{L^2}\right\rfloor,\qquad
\frac12\frac{n}{L^2}\le M\le \frac n{L^2},\qquad n\le M^2.
\tag{1}
\]
The nonempty injection support and normalization at this rounded size are
also supplied eventually by \(\noderef{TwoBiteNaturalTotalMassOneEventually}\).
Moreover
\[
pn=\frac12\sqrt{nL}\to\infty,\qquad \log M=O(L),
\tag{2}
\]
so \(pn/\log M\to\infty\).

Set
\[
\gamma=\min\{\eta/6,1/2\}.
\]
Then \(0<\gamma\le1/2\).  By
\(\noderef{FiberAndDegreeBaseDegreeConcentration}\), applied with tolerance
\(\gamma\), with high probability every red base degree and every blue base
degree is within relative error \(\gamma\) of \(pM\).  Denote this event by
\(\mathcal G_n\).  It is enough to show that, on \(\mathcal G_n\), the
conditional probability that some lifted neighborhood has size outside
relative error \(\eta\) from \(pn\) tends to \(0\), because
\(\Pr(\mathcal G_n)\to1\).

Fix a red base vertex \(r\in\operatorname{Fin}(M)\) and condition on the red
base graph \(G_R=\noderef{TwoBiteRedGraph}(\omega)\).  Let
\[
d_r=\big|\{u\in\operatorname{Fin}(M):G_R.Adj\,r\,u\}\big|.
\]
By the definition of \(\noderef{TwoBiteLiftedNeighborhood}\), together with
\(\noderef{TwoBiteRedProjection}\) and \(\noderef{TwoBiteEmbedding}\),
\(\noderef{TwoBiteLiftedNeighborhood}(\omega,\operatorname{inl}(r))\) is the
set of final vertices whose injected product cell has red coordinate in
\(N_{G_R}(r)\).  Equivalently, among the \(M^2\) product cells
\[
\operatorname{Fin}(M)\times\operatorname{Fin}(M),
\]
the marked cells are
\[
N_{G_R}(r)\times\operatorname{Fin}(M),
\]
and this marked set has cardinality \(d_rM\).

The two-bite sample weight factors by \(\noderef{TwoBiteSampleWeight}\) into
the red graph weight, the blue graph weight, and
\(\noderef{UniformInjectionWeight}\).  Conditional on \(G_R\), the event just
described depends only on the injection coordinate and not on the blue graph.
The blue graph weights sum to \(1\) by \(\noderef{GnpGraphWeightSumOne}\).
Since \(n\le M^2\), \(\noderef{UniformInjectionWeight}\) is the reciprocal of
the number of injections, so the injection coordinate is a uniform ordered
sample of \(n\) distinct cells from \(M^2\) cells.  Therefore, conditional on
\(G_R\), the random variable
\[
N_r=|\noderef{TwoBiteLiftedNeighborhood}(\omega,\operatorname{inl}(r))|
\]
is hypergeometric with population size \(M^2\), marked-set size \(d_rM\), and
sample size \(n\).  Its conditional mean is
\[
\mu_r=n\frac{d_rM}{M^2}=\frac{nd_r}{M}.
\tag{3}
\]

On \(\mathcal G_n\), the red degree satisfies
\[
|d_r-pM|\le \gamma pM.
\tag{4}
\]
Multiplying (4) by \(n/M\) gives
\[
|\mu_r-pn|\le \gamma pn.
\tag{5}
\]
Also \(\gamma\le1/2\) gives
\[
\mu_r\ge(1-\gamma)pn\ge \frac12 pn,
\qquad
\mu_r\le(1+\gamma)pn.
\tag{6}
\]

The node \(\noderef{HypergeometricMgfComparison}\) applies to the
hypergeometric variable \(N_r\).  Put \(\mu=\mu_r\) for this paragraph.  The
moment-generating-function comparison says that, for every real \(\lambda\),
\[
\mathbb E\!\left[e^{\lambda N_r}\mid G_R\right]
\le \exp\!\left(\mu(e^\lambda-1)\right).
\tag{7a}
\]
For the upper tail, take \(\lambda_+=\log(1+\gamma)>0\).  Markov's
inequality and (7a) give
\[
\begin{aligned}
\Pr(N_r\ge(1+\gamma)\mu\mid G_R)
&=\Pr(e^{\lambda_+N_r}\ge e^{\lambda_+(1+\gamma)\mu}\mid G_R)\\
&\le e^{-\lambda_+(1+\gamma)\mu}
   \mathbb E[e^{\lambda_+N_r}\mid G_R]\\
&\le
\exp\!\left(
  -\mu\bigl((1+\gamma)\log(1+\gamma)-\gamma\bigr)
\right).
\end{aligned}
\tag{7b}
\]
For the lower tail, take \(\lambda_-=-\log(1-\gamma)>0\) and apply (7a) with
\(-\lambda_-\).  On the event \(N_r\le(1-\gamma)\mu\), the random variable
\(e^{-\lambda_-N_r}\) is at least
\(e^{-\lambda_-(1-\gamma)\mu}\).  Thus Markov's inequality gives
\[
\begin{aligned}
\Pr(N_r\le(1-\gamma)\mu\mid G_R)
&\le e^{\lambda_-(1-\gamma)\mu}
   \mathbb E[e^{-\lambda_-N_r}\mid G_R]\\
&\le
\exp\!\left(
  \mu\bigl(-(1-\gamma)\log(1-\gamma)-\gamma\bigr)
\right)\\
&=
\exp\!\left(
  -\mu\bigl(\gamma+(1-\gamma)\log(1-\gamma)\bigr)
\right).
\end{aligned}
\tag{7c}
\]
For \(0<\gamma\le1/2\), both
\[
a_+=(1+\gamma)\log(1+\gamma)-\gamma,
\qquad
a_-=\gamma+(1-\gamma)\log(1-\gamma)
\]
are positive.  Define \(c_\gamma=\min\{a_+,a_-\}>0\).  Combining (7b) and
(7c) gives, conditionally on \(G_R\) and on \(\mathcal G_n\),
\[
\Pr\bigl(|N_r-\mu_r|>\gamma\mu_r\mid G_R\bigr)
\le 2\exp(-c_\gamma\mu_r)
\le 2\exp(-(c_\gamma/2)pn),
\tag{7}
\]
where the last inequality uses (6).

If the tail event in (7) does not occur, then using (5) and (6),
\[
|N_r-pn|
\le |N_r-\mu_r|+|\mu_r-pn|
\le \gamma\mu_r+\gamma pn
\le \gamma(1+\gamma)pn+\gamma pn
\le 3\gamma pn
\le \eta pn.
\tag{8}
\]
Thus the failure of
\(\noderef{WithinRelativeError}(\eta,pn,N_r)\) on \(\mathcal G_n\) is bounded
by the right side of (7).

The blue case is identical.  For a blue vertex \(b\), condition on
\(G_B=\noderef{TwoBiteBlueGraph}(\omega)\).  The marked cells are
\[
\operatorname{Fin}(M)\times N_{G_B}(b),
\]
whose size is \(d_bM\), and
\(\noderef{TwoBiteBlueProjection}\) identifies
\(\noderef{TwoBiteLiftedNeighborhood}(\omega,\operatorname{inr}(b))\) with
the injected vertices landing in those marked cells.  The same
hypergeometric estimate and the blue half of
\(\noderef{FiberAndDegreeBaseDegreeConcentration}\) give the same bound
\[
2\exp(-(c_\gamma/2)pn)
\]
for each fixed blue vertex.

There are \(M\) red and \(M\) blue base vertices.  By the union bound, the
conditional probability, on \(\mathcal G_n\), that any lifted neighborhood
violates the required relative-error bound is at most
\[
4M\exp(-(c_\gamma/2)pn).
\tag{9}
\]
By (2), the logarithm of (9) is
\[
\log 4+\log M-(c_\gamma/2)pn\to-\infty,
\]
so (9) tends to \(0\).  Combining this with
\(\Pr(\mathcal G_n)\to1\) proves that the simultaneous lifted-neighborhood
size concentration event has probability tending to \(1\), as required.
\end{proof}
